\scanpage{70}
Supposons $A$ \deleted{local} Dedekind \deleted{\unsure{$\ne$}} local \inserted{(cas Brauer)} de sorte que

\[
\begin{aligned}
\mathrm{Is}(C)&=1,\qquad \mathrm{Brd}(C)=\mathbb{Z}/\underset{\text{ramification}}{e}\mathbb{Z}\\
\mathrm{Pic}(A)&=0
\end{aligned}
\]

et on trouve

\[
\mathrm{Aut\,ext}(o)=\mathrm{Brd}\, C=\mathbb{Z}/e\mathbb{Z}
\]

Donc il y a obstruction locale à ce qu'un autom. de $o$ soit \unsure{néc.|nec.} intérieur, il y a $e$ classes d'autos de $o$, \ldots

Supposons \inserted{dans cas global} qu'un autom de $o$ soit localement intérieur (pour la top. \unsure{fppqc}, par exemple) \inserted{i.e. son image dans $\mathrm{Ram}(C\unsure{/\mathfrak p|/A})$ nulle}. \deleted{Alors les} \inserted{Il définit} donc une section de $\underline{o}^*/\mathrm{Gm}_A$, et \deleted{\unsure{elle est}} \inserted{l'obstruction} pour le relever est un élément de $\mathrm{Pic}(A)$, nul dans \deleted{\unsure{$\mathrm{Pic}(o)$}} $\mathrm{H}^1(A,o^*)$\ldots\ donnant un élément trivial de $\mathrm{Is}(C)$\ldots\ \deleted{\unsure{Donc}} ce qui se voyait aussi par les suites exactes précédentes

\begin{center}\rule{0.22\linewidth}{0.3pt}\end{center}

Cas des corps gauches de Deuring splittés en $p$ et $\infty$ : la suite exacte devient

\diagram{\begin{tikzcd}[ampersand replacement=\&,column sep=small]
0 \arrow[r,dash] \& \mathbb Z/2\mathbb Z\arrow[r] \& o^*\arrow[r] \& \mathrm{Aut}(o)\arrow[r] \& \mathbb Z/2\mathbb Z\arrow[r] \& \pi_0(C,x)\arrow[r] \& \pi_0(O,o)\arrow[r,dash] \&1
\end{tikzcd}}

\margin{engendré par $-1$}

\margin{$\Vert$ $\mathbb Z/2\mathbb Z$ sauf si $j=0$ ou $j=12^3$ (où on trouve $\mathbb Z/4\mathbb Z$ resp $\mathbb Z/6\mathbb Z$ si $p\ne 2,3$) \deleted{\unsure{et}}}

% slanted
\margin{on peut dire pour quels $p$ ces $j$ sont d'inv. de Hasse nul\ldots}

\margin{engendré par Frobenius}

\margin{épim ssi la courbe elliptique $X$ d'anneau $o$ se descend à $\mathbb F_p$, i.e. \unsure{ssi} $j\in\mathbb F_p$}

\margin{classes d'iso. des classes de courbes elliptiques de Hasse nul l'opération de ``Frobenius'' dessus est $E\longmapsto E^{(p)}$}

\margin{classes de conjugaison possibles d'ordres maximaux}

Question : $o^*$ est-il tjrs égal
% archive sheet number: 69

